``
`

\begin{itemize}
    \item
    \bigbreak
    \item SgAgr = Default Agree: When Merge applies first, both \(\uparrow \text{Agree} \uparrow\) and \(\downarrow \text{Agree} \downarrow\) succeed, leading to a clash, which is repaired by Default agreement
    \bigbreak
    \item PlAgr = NAgr, due to \ref{ex:PlAgr-Stojkovic}: 
    \begin{itemize}
        \item \(\uparrow \text{Agree} \uparrow\) fails, Merge applies, and then K successfully \(\downarrow \text{Agree} \downarrow\) with the features of N
        \item The verb agrees with the KP specified in this way (plural, feminine)
    \end{itemize}
\end{itemize}

\ea
\gll Sunčica nas zatim vodi do kuće u kojoj živi sa roditeljima. Njenih pet sestara živi u Pirotu, a brat je na školovanju.\\
      Sunčica us.\textsc{acc} then leads to house in which lives with parents.\textsc{instr} her.\textsc{gen} five sisters.\textsc{gen} live.\textsc{3sg} in Pirot and brother is at schooling.\textsc{loc}\\
\glt `Sunčica then takes us to the house where she lives with her parents. Her five sisters live in Pirot, and her brother is in school.'
\z

\ea 
\gll Osuđenici iz Niša su u tuči zadobili teže povrede. Posle ukazane medicinske pomoći oni su vraćeni u zatvor. \textit{Svih} \textit{šest} \textit{zatvorenika} nalaze se u Odeljenju pojačanog nadzora.\\
      convicts from Niš \textsc{aux}.3\textsc{pl} in fight sustained serious injuries after given medical assistance they \textsc{aux}.3\textsc{pl} returned to prison all six prisoners locate \textsc{refl} in Department increased surveillance\\
\glt `The convicts from Niš sustained serious injuries in a fight. After receiving medical assistance, they were returned to prison. All six prisoners are in the Enhanced Surveillance Department.'
\z

\ea 
\gll Derbi između Beka i Čačana razočarao je ljubitelje šaha. Naime, kapiteni su se dogovorili da se svih šest partija završi podelom poena, bez borbe.\\
      derby between Beka and Čačani disappointed \textsc{aux}.3\textsc{sg} fans chess.\textsc{gen} namely captains \textsc{aux}.3\textsc{pl} \textsc{refl} agreed that \textsc{refl} all six games finish.\textsc{3sg} split points without fight\\
\glt `The derby between Beka and Čačani disappointed chess fans. Namely, the captains agreed that all six games would end in a draw, without a fight.'
\z

\ea 
\gll Nakon što je odgovorni direktor Serđo Markjone prethodne nedelje najavio da šest fabrika u Italiji rizikuju zatvaranje, napetost između torinske automobilske kuće i radničkih sindikata sve je veća.\\
      after that \textsc{aux}.3\textsc{sg} responsible director Sergio Marchionne previous week announced that six factories in Italy risk.\textsc{3pl} closure tension between Turin automotive house and workers’ unions all is greater\\
\glt `After the responsible director Sergio Marchionne announced last week that six factories in Italy are at risk of closure, the tension between the Turin automotive company and the workers’ unions is growing.'
\z

Solution?
\begin{itemize}
    \item As in my proposal, SgAgr is agreement with a Sg situational \textit{pro}, which functions as the subject of thetic-like sentences, as they \textit{must} be anchored via a situation pronoun (on such \textit{pro}, see also \citealt{Erteschik-Shir1997, Basilico2003, Hinterholzl2019, Hinterholzl2024subjects, Hinterholzl2024es}) 
    \item The fact that SgAgr is also possible with definite subjects suggests that in Finnish, like in Serbian, referential numeral subjects \textit{can} stay low if the sentence is intended to be about a situation (thus again agreeing with \textit{pro}); for low referential subjects in sentences anchored via a situation pronoun, see also \citet{Hinterholzl2019, Hinterholzl2024es, Hinterholzl2024subjects}.
\end{itemize}

\subsection{{Definiteness and agreement cross-linguistically}}

\begin{itemize}
    \item Definite/specific NumPs in Russian obligatorily give rise to PlAgr \citep{Pereltsvaig2006, Matushansky2014, Matushansky2015a, Matushansky2015}
    \begin{itemize}
        \item \citet{Pereltsvaig2006}: PlAgr with DP, specified for $\phi$-features, SgAgr with "Small Nominals", without the D-layer.
        \item \citet{ Matushansky2014, Matushansky2015a, Matushansky2015}: SgAgr due to the lack of the individuation feature, responsible for specificity, i.e. D-linking, existence presupposition, referentiality, etc.
        \item But see \citet{Boskovic2008}: PlAgr due to agreement with the nominative subjects, SgAgr in the absence of the nominative.
    \end{itemize}
\end{itemize}

\ex.
\ag. Eti pjat' devušek \textcolor{red}{rabotali} / \textcolor{blue}{*rabotalo} tam.\\
     these five girls.\textsc{gen} worked.\textsc{n.sg} / *worked.\textsc{n.sg} there\\
    \glt `These five girls worked there.'\\
\bg. Pjat' etih devušek \textcolor{red}{rabotali} / \textcolor{blue}{rabotalo} tam.\\
     five these.\textsc{gen} girls.\textsc{gen} worked.\textsc{pl} / worked.\textsc{n.sg} there\\
    \glt `Five of these girls worked there.'

\begin{itemize}
    \item Similarly in Kazym Khanty: Definite NumPs obligatorily PlAgr \citep{Pisarenko2024}; \citet{Pisarenko2024} follows \citet{Pereltsvaig2006}'s analysis.
\end{itemize}

\begin{itemize}
    \item In Finnish, "semantic" definiteness (maximality, familiarity, distributivity) enables PlAgr with NumPs (even though the numeral itself is morpho-syntactically marked for singular) \citep{Kaiser2022salt, Kaiser2022short}.
    \begin{itemize}
        \item See also \citet{Norris2021}, who additionally discusses Estonian.
    \end{itemize}
    \item E.g., PlAgr in \ref{ex:Fin-pl} is compatible with a maximal interpretation (there are four parties in the context). SgAgr in \ref{ex:Fin-sg} can also be interpreted in an existential, non-maximal way: there exist four parties who reached an agreement.
    \item Importantly, SgAgr -- at least for some speakers -- is compatible with both maximal and non-maximal readings.
\end{itemize}

\ex.
\gll Neljä osapuol-ta \textcolor{red}{pää-sivät} sopimuks-een.\\
     four.\textsc{sg} party-\textsc{sg.par} reach-\textsc{past.3pl} agreement-\textsc{ill}\\
\glt `(The) four parties reached an agreement.' \label{ex:Fin-pl}

\ex.
\gll Neljä osapuol-ta \textcolor{blue}{pää-si} sopimuks-een.\\
     four.\textsc{sg} party-\textsc{sg.par} reach-\textsc{past.3sg} agreement-\textsc{ill}\\
\glt `Four parties reached an agreement.' / `There exist four parties that reached an agreement.' \label{ex:Fin-sg}

\citet{Kaiser2022salt}'s proposal
\begin{itemize}
    \item PlAgr is due to the definiteness/maximality effect.
    \item SgAgr is a consequence of the existential force and is actually \textit{absence of agreement}.
    \begin{itemize}
        \item Existential subjects presumably low in the structure.
        \item Existential construal = thetic sentence.
    \end{itemize}
    \bigbreak
\end{itemize}
Problem:
\begin{itemize}
    \item What about SgAgr with definite subject NumPs?
\end{itemize}

Solution?
\begin{itemize}
    \item As in my proposal, SgAgr is agreement with a Sg situational \textit{pro}, which functions as the subject of thetic-like sentences, as they \textit{must} be anchored via a situation pronoun (on such \textit{pro}, see also \citealt{Erteschik-Shir1997, Basilico2003, Hinterholzl2019, Hinterholzl2024subjects, Hinterholzl2024es}) 
    \item The fact that SgAgr is also possible with definite subjects suggests that in Finnish, like in Serbian, referential numeral subjects \textit{can} stay low if the sentence is intended to be about a situation (thus again agreeing with \textit{pro}); for low referential subjects in sentences anchored via a situation pronoun, see also \citet{Hinterholzl2019, Hinterholzl2024es, Hinterholzl2024subjects}.
\end{itemize}

Beyond 5-and-ups

\begin{itemize}
    \item PlAgr with "pluringulars" are often judged as ungrammatical in indefinite/nonspesific contexts \citep[106]{Danon2014}:

\ex.
\a. The/my family \textbf{are} making important decisions.
\b. \textsuperscript{??/*}One family \textbf{are} making important decisions.

\item \citet{Sauerland2004snippets}: PlAgr due to a null definite.

\bigbreak
\item D mediating verbal agreement with coordinated subjects \citep{mitic-arsenijevic-2019}
\item Definite null pronoun mediating relative pronouns' agreement in non-restrictive relative clauses (\citealt{Arsenijevic2016pluringulars}; building on \citealt{Postal1994, deVries2002, deVries2006}).

\bigbreak

\item D and clitic doubling; "semantic" agreement with pronouns (see \citealt{Milosavljevic2017} for an overview and discussion).

\end{itemize}


\subsubsection{Acceptability data} \label{subsect:acceptabiulity-data}

We now turn to the acceptability data based on \citet{Milosavljevic2018nasledje}. These are based on NumPs with numerals \textit{pet} 'five' and \textit{šest} 'six', with type of agreement (SgAgr vs. PlAgr), animacy (animate vs. inanimate) and agentivity (agentive vs. non-agentive) as independent variables, and the acceptability on the scale 1-5 as dependent variable (where 1 is completely unacceptable, and 5 completely acceptable). 384 native speakers participated in the experiment (see \citealt{Milosavljevic2018nasledje} for details). Table \ref{tab:animacy-agentivity} summarizes just the mean rates. PlAgr is significantly more acceptable across the board; see \citet{Milosavljevic2018nasledje} for a detailed statistical analysis.


\begin{table}[h]
    \centering
    \caption{Mean acceptability ratings of agreement patterns across animacy and agentivity categories}
    \label{tab:animacy-agentivity}
    \begin{tabular}{|c | c c c |} 
        \hline
        \textsc{Agr Type} & \textsc{Anim, Agent} & \textsc{Anim, Non-agent} & \textsc{Inanim, Non-agent} \\ 
        \hline\hline
        \textsc{SgAgr} & 3.45 & 3.20 & 3.22 \\ 
        \hline
        \textsc{PlAgr}  & 3.90 & 4.10 & 4.03 \\ 
        \hline
    \end{tabular}
\end{table}


\noindent Examples (\ref{ex:anim-agent}), (\ref{ex:anim-non-agent}) and (\ref{ex:non-anim-non-agent}) illustrate different combinations with respect to agentivity and animacy (in all of them definiteness/specificity was forced by possessive adjectives left to the numeral). 


\begin{itemize}
    \item Animate agentive subjects
\end{itemize}

\ea \label{ex:anim-agent}
\gll Mojih pet drugova popravlja / popravljaju auto u garaži. \\
    my.\textsc{gen.pl} five friends.\textsc{gen.pl} fix.\textsc{sg} / fix.\textsc{pl} car.\textsc{acc.sg} in garage.\textsc{loc.sg} \\
\glt ‘My five friends are fixing the car in the garage.’
\z


\begin{itemize}
    \item Animate non-agentive subjects
\end{itemize}

\ea \label{ex:anim-non-agent}
\gll Mirinih šest saradnika izgleda / izgledaju kao dobar tim. \\
    Mira's.\textsc{gen.pl} six colleagues.\textsc{gen.pl} seem.\textsc{sg} / seem.\textsc{pl} like good team.\textsc{nom.sg} \\
\glt ‘Mira’s six colleagues seem like a good team.’
\z

    

\begin{itemize} 
    \item Inanimate non-agentive subjects
\end{itemize}

\ea \label{ex:non-anim-non-agent}
\gll Mojih šest kofera teži / teže ukupno dvadeset kilograma. \\
    my.\textsc{gen.pl} six suitcases.\textsc{gen.pl} weigh.\textsc{sg} / weigh.\textsc{pl} total.\textsc{acc.sg} twenty kilograms.\textsc{gen.pl} \\
\glt ‘My six suitcases weigh a total of twenty kilograms.’

\footnote{\label{fn:double-case}According to \citet{Boskovic2008}, the numeral subject can be assigned case twice: in the base-generated position (genitive), and after moving to Spec,IP (nominative). The "double" case morphologically surfaces as genitive in Serbo-Croatian (though in Russian, it is realized as Nominative for high, Spec,IP subjects). It follows that low referential subjects (triggering SgAgr) are not assigned nominative, but only genitive. Consequently, only referential subjects that move to Spec,IP can be assigned nominative and can thus trigger PlAgr.}